\chapter{A Glimpse Beyond: Quaternions and Rotations}

\begin{center}\itshape
Hold a plate flat on your palm and turn it one full turn, keeping it horizontal: your arm is twisted.\\
Keep turning in the same direction: after a second full turn, your arm is straight again.\\
Why two turns, and not one?
\end{center}

Take two quarter-turns of an object in space, one about the horizontal axis and one about the vertical axis. Reverse their order. The final orientation may change. Three-dimensional rotation is not commutative, and any algebra meant to describe it must remember that.

Complex multiplication describes rotations in a plane. To describe rotations in space, we need four real coordinates rather than two. The \emph{quaternions} are numbers of the form
\[
q=a+bi+cj+dk,\qquad i^2=j^2=k^2=ijk=-1,
\]
with multiplication that satisfies $ij=k$ but $ji=-k$. Their lack of commutativity is not a defect. It is what makes the algebra fit the geometry.

A vector in $\R^3$ can be written as the imaginary quaternion $v=xi+yj+zk$. If $q$ is a unit quaternion, the map $v\mapsto qvq^{-1}$ is a rotation. The same rotation comes from $-q$. This is the double cover of the rotation group by the unit quaternions.

The groups we studied earlier return in a visible form. The rotations of the tetrahedron form $A_4$; those of the cube form $\mathfrak S_4$; and those of the icosahedron form $A_5$. In fact every finite subgroup of the three-dimensional rotation group belongs to one of five families.

Before closing, we will also multiply quaternion norms and recover the four-square identity. It is a good final example of how a useful algebraic rule can make a numerical identity seem inevitable.

\section{Quaternions}

Instead of postulating Hamilton's rules and checking associativity by hand, we realise quaternions as complex matrices, for which associativity is automatic.

\begin{definition}
Let $\mathbb H$ be the set of complex matrices
\[
\begin{pmatrix}z&w\\-\bar w&\bar z\end{pmatrix},\qquad z,w\in\C.
\]
It is a real vector space of dimension $4$, with basis $1=I$ and
\[
i=\begin{pmatrix}i&0\\0&-i\end{pmatrix},\quad j=\begin{pmatrix}0&1\\-1&0\end{pmatrix},\quad k=\begin{pmatrix}0&i\\i&0\end{pmatrix};
\]
so $q=a+bi+cj+dk$ corresponds to $z=a+bi$, $w=c+di$ (here $i$ on the right is the complex number). Its elements are \emph{quaternions}. The \emph{conjugate} of $q$ is $\bar q=a-bi-cj-dk$, its \emph{norm} is $N(q)=a^2+b^2+c^2+d^2$, its \emph{real part} is $a$, and $q$ is \emph{pure} if $a=0$; the pure quaternions form a copy of $\R^3$.
\end{definition}

\begin{proposition}\label{prop:quat}
\begin{enumerate}[label=(\alph*)]
\item $\mathbb H$ is closed under matrix multiplication; it is an associative ring, and
\[
\begin{gathered}i^2=j^2=k^2=-1,\qq ij=-ji=k,\quad jk=-kj=i,\qq ki=-ik=j.\end{gathered}
\]
\item $\bar q$ is the conjugate transpose of the matrix $q$, so $\overline{pq}=\bar q\,\bar p$; and $q\bar q=\bar qq=N(q)=\det q$. Hence $N(pq)=N(p)N(q)$.
\item Every $q\ne0$ has an inverse $q^{-1}=\bar q/N(q)$: $\mathbb H$ is a \emph{division ring} (a non-commutative field).
\item For pure quaternions $u,v$, identified with vectors of $\R^3$, $uv=-u\cdot v+u\times v$. In particular $u^2=-|u|^2$, and $u,v$ anticommute iff they are orthogonal.
\end{enumerate}
\end{proposition}
\begin{proof}
(a) The product of two such matrices is
\[
\begin{gathered}\begin{pmatrix}z&w\\-\bar w&\bar z\end{pmatrix}\begin{pmatrix}z'&w'\\-\bar w'&\bar z'\end{pmatrix}\qq=\begin{pmatrix}zz'-w\bar w'&zw'+w\bar z'\\-\bar wz'-\bar z\bar w'&-\bar ww'+\bar z\bar z'\end{pmatrix},\end{gathered}
\]
whose second row is $(-\bar W,\bar Z)$ for the first row $(Z,W)$. Matrix multiplication is associative and distributive. The relations are $2\times2$ products; for instance $ij=\begin{psmallmatrix}0&i\\i&0\end{psmallmatrix}=k$, and $ji=-k$.

(b) The conjugate transpose of $\begin{psmallmatrix}z&w\\-\bar w&\bar z\end{psmallmatrix}$ is $\begin{psmallmatrix}\bar z&-w\\\bar w&z\end{psmallmatrix}$, the matrix of $\bar z,-w$, i.e.\ of $a-bi-cj-dk$. The conjugate transpose reverses products. The matrix $q\bar q$ has entries $|z|^2+|w|^2$ on the diagonal and $0$ elsewhere, as does $\bar qq$; and $\det q=|z|^2+|w|^2=N(q)$. The determinant is multiplicative.

(c) $q\cdot\bar q/N(q)=\bar q/N(q)\cdot q=1$ by (b).

(d) Both sides are bilinear in $(u,v)$, so it suffices to check basis vectors: $i\cdot i=-1=-i\cdot i+i\times i$, and $ij=k=-0+i\times j$, and similarly for the other pairs. If $u\perp v$ then $vu=v\times u=-u\times v=-uv$; conversely $uv+vu=-2u\cdot v$.
\end{proof}

Taking norms in $(a+bi+cj+dk)(a'+b'i+c'j+d'k)$ gives Euler's four-square identity (Exercise~\ref{ex:euler4}), used in Volume~2 to reduce the four-square theorem to primes. Frobenius proved in 1878 that $\R$, $\C$ and $\mathbb H$ are the only finite-dimensional associative division algebras over $\R$; dropping associativity adds Cayley's octonions, of dimension $8$, and nothing else.

A non-zero quaternion has an inverse, and its norm behaves multiplicatively. For a quaternion of norm one, conjugation preserves the lengths of imaginary quaternions. We can now ask which rotation of space this conjugation describes.

\section{Unit quaternions and rotations}

The unit quaternions form a group $S^3=\{q:N(q)=1\}$ under multiplication (by Proposition~\ref{prop:quat}(b)), geometrically the unit sphere of $\R^4$. As matrices they are exactly $SU(2)$, the unitary $2\times2$ matrices of determinant $1$: such a matrix $\begin{psmallmatrix}\alpha&\beta\\\gamma&\delta\end{psmallmatrix}$ has inverse $\begin{psmallmatrix}\delta&-\beta\\-\gamma&\alpha\end{psmallmatrix}$ (determinant $1$) equal to its conjugate transpose, so $\delta=\bar\alpha$ and $\gamma=-\bar\beta$.

\begin{lemma}[Euler]\label{lem:euler}
Every $A\in SO(3)$ is a rotation about some axis: there are a unit vector $u$ and an angle $\theta$ such that $Au=u$ and $A$ turns the plane $u^\perp$ by $\theta$.
\end{lemma}
\begin{proof}
$\det(A-I)=\det(A-AA^{\mathsf T})=\det A\,\det(I-A^{\mathsf T})=\det(I-A)=-\det(A-I)$ (a $3\times3$ determinant changes sign with its matrix). So $\det(A-I)=0$ and $A$ has a unit vector $u$ with $Au=u$. Being orthogonal, $A$ preserves $u^\perp$; on this plane it is orthogonal with determinant $\det A=1$, hence a rotation by some $\theta$ (Volume~1).
\end{proof}

\begin{theorem}\label{thm:su2so3}
For a unit quaternion $q$ and a pure quaternion $v$, $R_q(v)=qvq^{-1}=qv\bar q$ is pure, and $R_q\in SO(3)$. The map $q\mapsto R_q$ is a surjective homomorphism $S^3\to SO(3)$ with kernel $\{\pm1\}$, so that
\[
SO(3)\cong S^3/\{\pm1\}=SU(2)/\{\pm I\}.
\]
Explicitly, if $u$ is a unit pure quaternion and $q=\cos\frac\theta2+\sin\frac\theta2\,u$, then $R_q$ is the rotation by $\theta$ about $u$.
\end{theorem}
\begin{proof}
\emph{First, conjugation preserves space and length.} A quaternion $v$ is pure exactly when $\bar v=-v$. If $v$ is pure and $q$ is a unit quaternion, then
\[
\overline{qv\bar q}=q\bar v\bar q=-qv\bar q.
\]
Thus $qv\bar q$ is pure as well. The map $R_q:v\mapsto qv\bar q$ is real-linear on the pure quaternions, which we identify with $\R^3$. The multiplicativity of the quaternion norm gives
\[
N(R_q(v))=N(q)N(v)N(\bar q)=N(v).
\]
Therefore $R_q$ is orthogonal. Moreover
\[
R_p(R_q(v))=p(qv\bar q)\bar p=(pq)v\overline{pq}=R_{pq}(v),
\]
so $q\mapsto R_q$ is a group homomorphism.

\emph{Find the rotation explicitly.} Write
\[
q=c+su,\qquad c=\cos(\theta/2),\quad s=\sin(\theta/2),
\]
where $u$ is a pure unit quaternion. Since $u$ commutes with $q$, we have $R_q(u)=u$. Now take a pure quaternion $w$ perpendicular to $u$. By Proposition~\ref{prop:quat}(d), $uw=u\times w=-wu$, and
\[
(u\times w)u=(u\times w)\times u=w
\]
because $u$ has unit length and is perpendicular to $w$. Expand:
\[
\begin{aligned}
qw\bar q
&=(cw+s\,u\times w)(c-su)\\
&=c^2w-cs\,wu+cs\,u\times w-s^2(u\times w)u\\
&=(c^2-s^2)w+2cs\,u\times w\\
&=\cos\theta\,w+\sin\theta\,u\times w.
\end{aligned}
\]
On the plane perpendicular to $u$, this is rotation through $\theta$, with its orientation fixed by $u$. The vector $u$ stays fixed. Hence $R_q\in SO(3)$.

\emph{Every rotation occurs.} By Lemma~\ref{lem:euler}, every element of $SO(3)$ is a rotation about some unit axis $u$ through some angle $\theta$. The formula above shows that it is $R_q$ for $q=\cos(\theta/2)+u\sin(\theta/2)$. So the map is onto.

\emph{Exactly two quaternions give each rotation.} If $R_q$ is the identity, then $q$ commutes with every pure quaternion. Write $q=a+x$, with $x$ pure. For every pure $w$,
\[
0=qw-wq=xw-wx=2x\times w.
\]
Only $x=0$ has zero cross product with \emph{every} $w$. Thus $q=a$ is real, and $N(q)=a^2=1$ gives $q=\pm1$. Conversely both $1$ and $-1$ act trivially. The kernel is therefore $\{1,-1\}$, and every rotation has exactly two preimages.
\end{proof}

\paragraph{Why this works.}
A unit quaternion $q$ preserves the quaternion norm under conjugation, because $N(qvq^{-1})=N(v)$. It also preserves the imaginary quaternions and therefore gives an orthogonal map of $\R^3$. The map depends continuously on $q$ and sends $1$ to the identity, so its determinant is $+1$, not $-1$: it is a rotation. A quaternion of the form $\cos(\theta/2)+u\sin(\theta/2)$ produces rotation through angle $\theta$ around the unit vector $u$. Both $q$ and $-q$ produce the same map; there are no other duplicates. The \emph{half-angle} is what makes the double cover work.

\begin{figure}[ht]
\centering
\fig{ch12-rotation}
\caption{The unit quaternion $q=\cos\frac\theta2+\sin\frac\theta2u$ turns every vector $v$ by the angle $\theta$ about the axis $u$: $v\mapsto qvq^{-1}$. The half-angle explains why $q$ and $-q$, which differ by a turn of $2\pi$ in $\theta$, give the same rotation.}
\label{fig:rotation}
\end{figure}

Composing rotations is multiplying quaternions (Exercise~\ref{ex:compose}): four numbers instead of nine, no drift away from orthogonality that cannot be cured by dividing by the norm, and smooth interpolation between orientations. This is why flight software, robots and video games store orientations as unit quaternions.

\begin{remark}[the plate trick]
Turning the plate by an angle $2\pi t$, $0\le t\le1$, about the vertical axis $k$ is the path $q(t)=\cos\pi t+\sin\pi t\,k$ in $S^3$, which goes from $1$ to $-1$: it does not close up, although the rotations $R_{q(t)}$ do. Two turns give $\cos2\pi t+\sin2\pi t\,k$, a closed loop on the $3$-sphere, which can be shrunk to a point on the sphere; its shrinking, projected to $SO(3)$, is the motion that untwists your arm. In topological language $S^3$ is simply connected and $SO(3)$ is not: its fundamental group is $\Z/2\Z$. In physics, electrons are described by vectors on which $S^3=SU(2)$, not $SO(3)$, acts, and a rotation by $2\pi$ multiplies them by $-1$: they have spin $\frac12$.
\end{remark}

\begin{think}
How many rotations map a regular icosahedron to itself? (Count where a given vertex can go, and how many rotations fix it.) Which group of Chapter~4 has that order and should be its rotation group?
\end{think}

Unit quaternions describe every rotation continuously. We now restrict attention to finite sets of rotations that are closed under composition. Their axes and angles can be counted, and the count leads back to the cyclic, dihedral and alternating groups.

\section{Finite groups of rotations}

For each regular polyhedron (Figure~\ref{fig:solids}), the rotations preserving it form a finite subgroup of $SO(3)$. A rotation is determined by the image of a vertex and of an adjacent vertex (a rotation fixing a vertex $x$ and an adjacent vertex $y$ fixes the two independent vectors $x,y$, hence their cross product, hence everything). For the icosahedron, a vertex can go to any of the $12$ vertices, and then its $5$ neighbours can be rotated in $5$ ways: $60$ rotations. In the same way the tetrahedron has $4\cdot3=12$ rotations, the cube and the octahedron $24$, the dodecahedron and the icosahedron $60$. The orbit--stabiliser theorem is behind each count.

\begin{figure}[ht]
\centering
\fig{ch12-solids}
\caption{The five Platonic solids and their rotation groups. Dual solids (cube and octahedron, dodecahedron and icosahedron: put a vertex at the centre of each face) have the same rotation group, and the tetrahedron is self-dual.}
\label{fig:solids}
\end{figure}

\begin{theorem}\label{thm:finiteso3}
A finite subgroup $G$ of $SO(3)$ is cyclic, or dihedral of order $2n\ge4$, or isomorphic to $A_4$, $\mathfrak S_4$ or $A_5$ (of orders $12$, $24$, $60$).
\end{theorem}
\begin{proof}
Let $N=|G|>1$. Each $g\ne1$ in $G$ is a rotation about an axis, which meets the unit sphere in two \emph{poles} $\pm u$, the only unit vectors fixed by $g$. Let $P$ be the set of poles of all non-identity elements. If $p$ is a pole of $g$ and $h\in G$, then $hp$ is a pole of $hgh^{-1}$: so $G$ acts on $P$. The stabiliser $G_p$ consists of the rotations of $G$ about the axis $\R p$; it is a finite group of rotations of the plane $p^\perp$, isomorphic to a finite subgroup of $\{z\in\C:|z|=1\}$, hence cyclic (Theorem~\ref{thm:cyclicmult}), of order $n_p\ge2$.

\emph{Counting.} Count the pairs $(g,p)$ with $g\ne1$ and $p$ a pole of $g$. Each $g\ne1$ has $2$ poles, and each pole $p$ belongs to $n_p-1$ elements $g\ne1$. Group the poles into orbits $O_1,\dots,O_r$; on $O_i$, $n_p=n_i$ is constant (conjugate stabilisers) and $|O_i|=N/n_i$ (orbit--stabiliser). So
\[
\begin{gathered}2(N-1)=\sum_{i=1}^r\frac N{n_i}(n_i-1),\qq\text{i.e.}\quad2-\frac2N=\sum_{i=1}^r\Bigl(1-\frac1{n_i}\Bigr).\end{gathered}
\]
The left side is in $[1,2)$, each term on the right in $[\frac12,1)$, so $r=2$ or $r=3$.

\emph{$r=2$.} Then $\frac2N=\frac1{n_1}+\frac1{n_2}$ with $n_i\le N$, so $n_1=n_2=N$: every pole is fixed by all of $G$. There are two of them, $\pm p$, and $G=G_p$ is cyclic.

\emph{$r=3$.} Order $n_1\le n_2\le n_3$. Then $1+\frac2N=\frac1{n_1}+\frac1{n_2}+\frac1{n_3}$, so $n_1=2$ (otherwise the right side is $\le1$), and $\frac12+\frac2N=\frac1{n_2}+\frac1{n_3}$. If $n_2\ge4$ the right side is $\le\frac12$; so $n_2\in\{2,3\}$. If $n_2=2$, then $n_3=N/2$. If $n_2=3$, then $\frac1{n_3}=\frac16+\frac2N>\frac16$ gives $n_3\in\{3,4,5\}$ and $N=12,24,60$. Four cases remain: $(n_1,n_2,n_3)=(2,2,n)$ with $N=2n$, and $(2,3,3)$, $(2,3,4)$, $(2,3,5)$.

\emph{$(2,2,n)$, $N=2n$.} If $n=2$, $N=4$ and every element has order $1$ or $2$, so $G\cong(\Z/2\Z)^2=D_2$. Let $n\ge3$. The orbit $O_3$ has $N/n=2$ poles; if $p\in O_3$, then $-p$ (with the same stabiliser) has $n_{-p}=n\ge3$, so $-p\in O_3$, and $O_3=\{\pm p\}$. Let $C=G_p$, cyclic of order $n$ and index $2$. An element $s\notin C$ maps $p$ to $-p$, so it is not a rotation about $\R p$. Its poles have stabiliser of order $2$ (otherwise they would be in $O_3$), so $s$ has order $2$: it is a half-turn about an axis perpendicular to $p$ (it reverses $p$). Conjugating by $s$ maps the rotation by $\alpha$ about $p$ to the rotation by $\alpha$ about $-p$, i.e.\ by $-\alpha$ about $p$: $scs^{-1}=c^{-1}$ for $c\in C$. With $C=\langle c\rangle$ of order $n$ and $G=C\cup sC$, this is the dihedral group $D_n$ (Chapter~1: the same multiplication rules).

\emph{$(2,3,3)$, $N=12$.} The orbit $O_2$ has $4$ poles, and $G$ permutes them. A non-identity rotation fixes only $2$ unit vectors, so only the identity fixes all four: $G$ embeds in $\mathfrak S_4$, as a subgroup of order $12$. A subgroup $H$ of index $2$ in $\mathfrak S_4$ is normal and contains the square of every element ($g^2H=(gH)^2=H$ in $\mathfrak S_4/H\cong\Z/2\Z$); every $3$-cycle is a square (of its inverse), so $H$ contains the $3$-cycles, which generate $A_4$ (Proposition~\ref{prop:gens}); so $H=A_4$ and $G\cong A_4$.

\emph{$(2,3,4)$, $N=24$.} $O_2$ has $8$ poles, with stabilisers of order $3$; it is the only orbit with $n=3$, so it contains $-p$ with $p$: it consists of $4$ lines $\ell_1,\dots,\ell_4$ through $0$, which $G$ permutes. Let $g\ne1$ fix every $\ell_m$. Let $r$ be the rotation of order $3$ about $\ell_1$ in $G$. A rotation by $\pm120^\circ$ maps no line to itself except its axis, so $r$ moves $\ell_2,\ell_3,\ell_4$. If $g$ had one of the $\ell_m$ as its axis, $g$ would be a rotation of order $3$ about it (it is in the stabiliser of a pole of $O_2$), and would move the other three lines: impossible. So no $\ell_m$ is the axis of $g$, and $g$ maps each $\ell_m$ to itself only by reversing it: $g$ is a half-turn and all $\ell_m$ are perpendicular to its axis, hence lie in one plane $\Pi$. But $r$ permutes the $\ell_m$, so it maps $\Pi$, the plane they span, to itself, and therefore also the line $\Pi^\perp$; as $r$ preserves only its axis $\ell_1$, $\Pi^\perp=\ell_1\subseteq\Pi$, a contradiction. So $G$ acts faithfully on $4$ lines and embeds in $\mathfrak S_4$; with $|G|=24$, $G\cong\mathfrak S_4$.

\emph{$(2,3,5)$, $N=60$.} The orbits have $30$, $20$ and $12$ poles, with stabilisers of orders $2$, $3$, $5$. Since an axis carries $2$ poles, with the same stabiliser, $G$ consists of the identity, $15$ half-turns, $10\cdot2=20$ elements of order $3$ and $6\cdot4=24$ of order $5$ (total $60$). First, $G$ is simple. Let $H\trianglelefteq G$, $H\ne\{1\}$, contain $g\ne1$, with pole $p$. As $G_p$ has prime order, $g$ generates it, so $G_p\subseteq H$, and then $G_{hp}=hG_ph^{-1}\subseteq H$ for all $h$. So $H$ contains every element whose poles lie in the orbit of $p$: $15$, $20$ or $24$ elements. Hence $|H|=1+$ a sum of distinct terms among $15,20,24$, and $|H|$ divides $60$; the only possibility is $1+15+20+24=60$ (the others are $16,21,25,36,40,45$). So $H=G$.

Next, the half-turns are all conjugate (their poles form one orbit), so each one, $a$, has a centraliser of order $60/15=4$. By Sylow's theorems (Theorem~\ref{thm:sylow23}), $a$ lies in a subgroup $P$ of order $4$, which is abelian (Corollary~\ref{cor:pgroup}), so $P\subseteq C_G(a)$ and $P=C_G(a)$. Elements of $P\setminus\{1\}$ have order $2$ ($G$ has no element of order $4$). So the Sylow $2$-subgroups each contain $3$ half-turns, each half-turn lies in exactly one of them, and there are $15/3=5$ of them. $G$ acts on these $5$ subgroups by conjugation, transitively (Sylow), so the homomorphism $G\to\mathfrak S_5$ is not trivial; its kernel is a normal subgroup other than $G$, hence trivial. So $G$ is isomorphic to a subgroup of index $2$ of $\mathfrak S_5$, which is $A_5$ by the argument used for $\mathfrak S_4$.
\end{proof}

\paragraph{Why this works.}
Every non-identity rotation fixes exactly two points of the unit sphere, the poles on its axis. Counting pairs (rotation, fixed pole) first by rotations gives $2(|G|-1)$. Counting first by pole orbits gives
\[
2-\frac2{|G|}=\sum_i\left(1-\frac1{n_i}\right),
\]
where $n_i$ is the order of a pole stabiliser. That short equation permits only two or three pole orbits. With three, the stabiliser orders must be $(2,2,n)$, $(2,3,3)$, $(2,3,4)$, or $(2,3,5)$. The rest of the proof recognises the corresponding group by how it permutes its axes. The numerical restriction comes from counting; identifying the abstract group needs the extra action arguments.

All these cases occur: rotations about one axis by multiples of $2\pi/n$, the symmetries of a regular $n$-gon in a plane (turning it over is a half-turn in space), and the rotation groups of the tetrahedron, the cube and the icosahedron, of orders $12$, $24$, $60$, which by the theorem are $A_4$, $\mathfrak S_4$ and $A_5$. One can show that two finite subgroups of $SO(3)$ that are isomorphic are conjugate; so up to a change of axes these are \emph{all} the finite rotation groups, and every finite symmetry of a molecule or a crystal, up to reflections, is in this list. In particular the rotation group of the icosahedron is $A_5$, and its natural $3$-dimensional representation is the character $\psi_3$ of Chapter~11.

The preimages of these groups in $S^3$ are the \emph{binary} polyhedral groups, of orders $24$, $48$ and $120$ (Exercises~\ref{ex:2T} and~\ref{ex:2I}). The binary icosahedral group $2I$, of order $120$, is especially rich: the quotient $S^3/2I$ is the Poincaré homology sphere, a $3$-manifold with the same homology as $S^3$ but a different fundamental group, and the $120$ elements of $2I$ are the vertices of the regular $4$-dimensional polytope called the $600$-cell. These connections, found one by one over a century, are still being explored.

The norm of a quaternion is a sum of four squares. Because norms multiply, multiplying two quaternions turns the product of two such sums into another sum of four squares. The resulting identity is a piece of number theory hidden in the algebra of rotations.

\section{Quaternions and four squares}

Volume~2 proved that every positive integer is a sum of four squares, by Lagrange's descent. Quaternions explain the theorem as unique factorisation explains the two-square theorem in $\Z[i]$. The \emph{Hurwitz quaternions} $\mathcal H$ are the quaternions $a+bi+cj+dk$ whose four coordinates $a,b,c,d$ are either all in $\Z$ or all in $\Z+\frac12$.
They form a ring in which division with remainder works on one side, and the proof that a prime is a norm follows the proof for $\Z[i]$; Exercise~\ref{ex:hurwitz} carries it out in full.

\begin{remark}[a first look at categories]
Many constructions of this volume share a pattern: they take objects (groups, rings, fields, representations) to other objects, and maps between objects to maps between the new objects, respecting composition. $G\mapsto G/G'$, $R\mapsto R[X]$, a finite abelian group $A\mapsto\widehat A$, $q\mapsto R_q$: the language of \emph{categories} (objects and morphisms) and \emph{functors} (maps between categories) was invented in 1945 by Eilenberg and Mac Lane to make such patterns precise. The duality $A\mapsto\widehat A$ reverses arrows (a homomorphism $A\to B$ gives $\widehat B\to\widehat A$), and the Galois correspondence of Chapter~10 is an order-reversing dictionary of the same kind. Volume~10 uses this language systematically.
\end{remark}

\begin{beauty}
Three facts from different chapters meet in the icosahedron: its rotation group is $A_5$, the smallest non-abelian simple group (Chapter~4), which is why the quintic has no formula (Chapter~10); its character table contains the golden ratio (Chapter~11); and its double cover in the quaternions is the binary icosahedral group, whose $120$ elements are the vertices of a regular polytope in four dimensions, the $600$-cell. Klein's \emph{Lectures on the Icosahedron} (1884) used these symmetries to solve the quintic, not by radicals, but with one extra function, exactly as the cubic is solved with trigonometry. Symmetry, algebra and arithmetic are one subject.
\end{beauty}

\section{Exercises}

\begin{exercise}[$\star$]\label{ex:compose}
Use quaternions to find the axis and angle of the rotation obtained by turning by $90^\circ$ about the $z$-axis, then by $90^\circ$ about the $x$-axis.
\end{exercise}

\begin{exercise}[$\star$]\label{ex:euler4}
Write Euler's identity expressing $(a^2+b^2+c^2+d^2)(a'^2+b'^2+c'^2+d'^2)$ as a sum of four squares of bilinear expressions, using $N(pq)=N(p)N(q)$.
\end{exercise}

\begin{exercise}[$\star\star$]\label{ex:rotmatrix}
Let $q=a+bi+cj+dk$ be a unit quaternion. Write the matrix of $R_q$ in the basis $i,j,k$, and check that its trace is $4a^2-1=1+2\cos\theta$, where $\theta$ is the angle of the rotation.
\end{exercise}

\begin{exercise}[$\star\star$]\label{ex:conjS3}
Show that two unit quaternions are conjugate in $S^3$ iff they have the same real part. Deduce that two rotations are conjugate in $SO(3)$ iff they have the same angle (in $[0,\pi]$).
\end{exercise}

\begin{exercise}[$\star\star$]\label{ex:2T}
Show that the $24$ quaternions $\pm1,\pm i,\pm j,\pm k,\frac12(\pm1\pm i\pm j\pm k)$ form a subgroup $2T$ of $S^3$, that its image in $SO(3)$ has order $12$, and that this image is isomorphic to $A_4$ (count its elements of order $2$).
\end{exercise}

\begin{exercise}[$\star\star$]\label{ex:2I}
Let $2I\subset S^3$ be the preimage of a rotation group $G\cong A_5$. Show that $|2I|=120$, that $-1$ is the only element of order $2$ in $S^3$, and deduce that $2I$ is isomorphic neither to $\mathfrak S_5$ nor to $A_5\times\Z/2\Z$.
\end{exercise}

\begin{exercise}[$\star\star\star$]\label{ex:hurwitz}
(Four squares with quaternions.)
\begin{enumerate}[label=(\alph*)]
\item Show that $\mathcal H$ is a subring of $\mathbb H$, stable under conjugation, and that $N(x)\in\Z$ for $x\in\mathcal H$. Show that the units of $\mathcal H$ (elements of norm $1$) are the $24$ elements of $2T$.
\item Show that for every real quaternion $x$ there is $h\in\mathcal H$ with $N(x-h)<1$. Deduce: for $a,b\in\mathcal H$, $b\ne0$, there are $q,r\in\mathcal H$ with $a=qb+r$ and $N(r)<N(b)$; and every left ideal of $\mathcal H$ is of the form $\mathcal Hb$.
\item Let $p$ be an odd prime. Using integers $x,y$ with $1+x^2+y^2\equiv0\pmod p$ (Volume~2, Chapter~7), show that $p=N(\delta)$ for some $\delta\in\mathcal H$, and then that $p$ is a sum of four integer squares.
\end{enumerate}
\end{exercise}

\begin{exercise}[$\star\star$]\label{pb:sylvester}
\begin{enumerate}[label=(\alph*)]
\item (Fibonacci, 1202.) Let $0<\frac ab<1$ be rational. Subtract the largest unit fraction $\frac1n\le\frac ab$, and repeat
with the remainder. Show that the process stops, and writes $\frac ab$ as a sum of distinct unit fractions. Apply it to
$\frac4{13}$.
\item (Sylvester, 1880.) Let $s_0=2$ and $s_{k+1}=s_k^2-s_k+1$, so $2$, $3$, $7$, $43$, $1807$, \dots Show that
$\sum_{k=0}^N\frac1{s_k}=1-\frac1{s_{N+1}-1}$, and that the $s_k$ are pairwise coprime.
\item Show that there is a constant $E>1$ such that $s_k=\bigl\lfloor E^{2^{k+1}}+\frac12\bigr\rfloor$ for every $k\ge0$.
\end{enumerate}
\end{exercise}

\begin{exercise}[$\star\star$]\label{pb:jensen}
\begin{enumerate}[label=(\alph*)]
\item Show that for real $x$ and $\theta$ and $n\ge1$,
\[
\prod_{k=0}^{n-1}\Bigl(x^2-2x\cos\Bigl(\theta+\frac{2\pi k}n\Bigr)+1\Bigr)=x^{2n}-2x^n\cos n\theta+1.
\]
\item Deduce Poisson's integral: for $|x|\ne1$, $\int_0^\pi\log(1-2x\cos\theta+x^2)\dd\theta$ is $0$ if $|x|<1$ and
$2\pi\log|x|$ if $|x|>1$.
\item Let $P=a\prod_{k=1}^d(X-z_k)\in\C[X]$, $a\ne0$, with no root of modulus $1$. Show that
\[
\frac1{2\pi}\int_0^{2\pi}\log\bigl|P(e^{i\theta})\bigr|\dd\theta=\log|a|+\sum_{k=1}^d\log\max(1,|z_k|).
\]
What does Problem~\ref{pb:kronecker} say about the polynomials of $\Z[X]$ for which this is $0$?
\end{enumerate}
\end{exercise}

\begin{exercise}[$\star$]\label{pb:mersennediv}
Let $p$ be an odd prime and $M_p=2^p-1$.
\begin{enumerate}[label=(\alph*)]
\item Show that every prime factor $q$ of $M_p$ satisfies $q\equiv1\pmod{2p}$ and $q\equiv\pm1\pmod8$.
\item Show that if $p\equiv3\pmod4$ and $2p+1$ is prime, then $2p+1$ divides $M_p$. Deduce that $M_{11}$, $M_{23}$ and
$M_{83}$ are composite.
\end{enumerate}
\end{exercise}

\begin{exercise}[$\star$]\label{pb:dice}
A die is a cube whose faces carry the numbers $1$ to $6$, each once. Two dice are the same if a rotation carries one onto the
other.
\begin{enumerate}[label=(\alph*)]
\item How many different dice are there?
\item On a standard die, opposite faces add up to $7$. How many different standard dice are there?
\item Show that no die is the same as its mirror image, so that the dice of (a) come in mirror pairs.
\end{enumerate}
\end{exercise}

\begin{exercise}[$\star\star\star$]\label{pb:ksubsets}
Let $\pi_k(\sigma)$ be the number of $k$-element subsets of $\{1,\dots,n\}$ fixed by $\sigma\in\mathfrak S_n$, and let
$0\le k,l\le\frac n2$.
\begin{enumerate}[label=(\alph*)]
\item Show that $\langle\pi_k,\pi_l\rangle=\min(k,l)+1$.
\item Deduce that $\chi_0=1$ and $\chi_k=\pi_k-\pi_{k-1}$ ($1\le k\le\frac n2$) are distinct irreducible characters of
$\mathfrak S_n$, of degrees $\binom nk-\binom n{k-1}$.
\item (Livingstone and Wagner.) Let $G\le\mathfrak S_n$. Show that the number of orbits of $G$ on $k$-subsets is a
non-decreasing function of $k$ for $0\le k\le\frac n2$.
\end{enumerate}
\end{exercise}

\begin{exercise}[$\star\star$]\label{pb:egz}
\emph{Erdős--Ginzburg--Ziv (1961).} We show that among any $2n-1$ integers there are $n$ whose sum is divisible by $n$.
\begin{enumerate}[label=(\alph*)]
\item Prove Chevalley--Warning for systems: if $f_1,\dots,f_r\in\mathbb F_p[x_1,\dots,x_m]$ and $\sum_j\deg f_j<m$, then the number
of common zeros of the $f_j$ in $\mathbb F_p^m$ is divisible by $p$.
\item Prove the theorem for $n=p$ prime, with $f_1=\sum_ia_ix_i^{p-1}$ and $f_2=\sum_ix_i^{p-1}$.
\item Show that if the theorem holds for $m$ and for $n$, it holds for $mn$. Conclude.
\item Show that $2n-1$ cannot be replaced by $2n-2$.
\end{enumerate}
\end{exercise}

\begin{exercise}[$\star\star$]\label{pb:abelian}
\emph{How many abelian groups? (Erdős and Szekeres, 1934).} Let $A(n)$ be the number of abelian groups of order $n$, up to
isomorphism.
\begin{enumerate}[label=(\alph*)]
\item Show that $A(n)=\prod_ip(e_i)$ if $n=\prod_ip_i^{e_i}$, and that $\sum_nA(n)n^{-s}=\prod_{k\ge1}\zeta(ks)$ for $s>1$.
\item Show that $\sum_{n\le x}A(n)=cx+O(\sqrt x)$ with $c=\prod_{k\ge2}\zeta(k)=2.2948\ldots$: on average, there are about $2.29$
abelian groups of each order.
\end{enumerate}
\end{exercise}

\begin{exercise}[$\star\star\star$]\label{pb:dodecaturns}
\emph{Random turns of an icosahedron.} The rotation group of the regular icosahedron is isomorphic to $A_5$ (Theorem~12.5), and the
$20$ rotations by a third of a turn about the axes through the centres of opposite faces form the conjugacy class of the
$3$-cycles. An icosahedron is turned, again and again, by a uniformly chosen one of these $20$ rotations.
\begin{enumerate}[label=(\alph*)]
\item Let $\mu$ be a probability on a finite group $G$ that is constant on conjugacy classes, and $\mu^{*t}$ the law of the walk at time
$t$, started at $e$. Show that, for every $g\in G$,
\[
\mu^{*t}(g)=\frac1{|G|}\sum_\chi\chi(1)\overline{\chi(g)}\,\Bigl(\frac{\hat\mu(\chi)}{\chi(1)}\Bigr)^t,\qquad\hat\mu(\chi)=\sum_{h\in G}\mu(h)\chi(h),
\]
the sum being over the irreducible characters of $G$. (Show first that, for a class function $f$ and an irreducible representation $\rho$ of
dimension $d$, the matrix $\sum_hf(h)\rho(h)$ is the scalar $\frac1d\sum_hf(h)\chi_\rho(h)$.)
\item Using the character table of $A_5$ (Exercise~11.7), show that the icosahedron is back in its initial position after $t$
turns with probability
\[
\frac1{60}\Bigl(1+16\cdot4^{-t}+25\cdot\Bigl(-\frac15\Bigr)^t\Bigr).
\]
\item Find the probability that after $t$ turns the icosahedron has turned by a fifth of a turn (in either direction) about some
axis through opposite vertices, and show that the law of the position converges to the uniform law at the rate $4^{-t}$.
\end{enumerate}
\end{exercise}

\section{Hints}

\begin{hint}{ex:compose}
The quaternions are $\frac{1+k}{\sqrt2}$ and $\frac{1+i}{\sqrt2}$; the composite ``first $q_1$, then $q_2$'' is $R_{q_2q_1}$.
\end{hint}
\begin{hint}{ex:euler4}
Multiply $p=a+bi+cj+dk$ by $q=a'+b'i+c'j+d'k$ with the rules of Proposition~\ref{prop:quat}.
\end{hint}
\begin{hint}{ex:rotmatrix}
Compute $qi\bar q$, $qj\bar q$, $qk\bar q$, or use Theorem~\ref{thm:su2so3}: the trace of a rotation by $\theta$ is $1+2\cos\theta$.
\end{hint}
\begin{hint}{ex:conjS3}
$q=\cos\alpha+\sin\alpha\,u$, and $pup^{-1}$ can be any unit pure quaternion.
\end{hint}
\begin{hint}{ex:2T}
Check closure on the generators $i$ and $\omega=\frac12(1+i+j+k)$, or check that the $24$ elements are the units of $\mathcal H$. Real parts: $\pm1$, $0$, $\pm\frac12$.
\end{hint}
\begin{hint}{ex:2I}
The kernel $\{\pm1\}$ lies in $2I$. Solve $q^2=1$ with $q=a+x$.
\end{hint}
\begin{hint}{ex:hurwitz}
(b) Round each coordinate; the bad case is when all four errors are $\frac12$. Then take $q$ close to $ab^{-1}$. (c) Look at the left ideal $\mathcal Hp+\mathcal H(1+xi+yj)$.
\end{hint}

\begin{hint}{pb:sylvester}
(b) $s_{k+1}-1=s_k(s_k-1)$. (c) Put $u_k=s_k-\frac12$; then $u_{k+1}=u_k^2+\frac14$. Study $2^{-k}\log u_k$.
\end{hint}

\begin{hint}{pb:jensen}
(a) $x^2-2x\cos\varphi+1=(x-e^{i\varphi})(x-e^{-i\varphi})$. (b) Riemann sums at the points $\frac{2\pi k}n$. (c) Apply (b)
to each factor.
\end{hint}

\begin{hint}{pb:mersennediv}
(a) The order of $2$ modulo $q$ is $p$, and $2=2^{p+1}$ is a square modulo $q$. (b) $2p+1\equiv7\pmod8$.
\end{hint}

\begin{hint}{pb:dice}
The $24$ rotations of the cube act on numberings. Which rotation fixes a numbering?
\end{hint}

\begin{hint}{pb:ksubsets}
(a) The orbits of $\mathfrak S_n$ on pairs $(A,B)$ of subsets are determined by $|A\cap B|$. (b) A virtual character of norm
$1$ and positive degree is irreducible. (c) The number of orbits of $G$ is $\langle\pi_k,1\rangle_G$.
\end{hint}

\begin{hint}{pb:egz}
(a) The number of common zeros is $\sum_x\prod_j(1-f_j(x)^{p-1})$ modulo $p$. (c) Extract groups of $n$ one at a time.
\end{hint}

\begin{hint}{pb:abelian}
(a) Corollary~7.7. (b) $A=\mathbf1*g$ with $g\ge0$ supported on powerful numbers.
\end{hint}

\begin{hint}{pb:dodecaturns}
(a) $\sum_hf(h)\rho(h)$ commutes with every $\rho(g)$; Schur's lemma, then the orthogonality relations to expand the class function $\mu^{*t}$ in
characters. (b) The ratios $\frac{\chi(g)}{\chi(1)}$ on a $3$-cycle are $1,0,0,\frac14,-\frac15$. (c) On a $5$-cycle the characters of dimensions $4$ and $5$ take the values
$-1$ and $0$.
\end{hint}

\section{Solutions}

\begin{solution}{ex:compose}
\idea Multiply the two unit quaternions in the right order and read off the half-angle and the axis.
\writeup $q_1=\cos45^\circ+\sin45^\circ\,k=\frac{1+k}{\sqrt2}$ and $q_2=\frac{1+i}{\sqrt2}$. The composite $R_{q_2}R_{q_1}=R_{q_2q_1}$, and
\[
\begin{aligned}q_2q_1&=\tfrac12(1+i)(1+k)\eqbr=\tfrac12(1+k+i+ik)\eqbr=\tfrac12(1+i-j+k),\end{aligned}
\]
since $ik=-j$. Its real part is $\cos\frac\theta2=\frac12$, so $\theta=120^\circ$, and its pure part $\frac12(i-j+k)=\sin60^\circ\cdot\frac{i-j+k}{\sqrt3}$ gives the axis $\frac1{\sqrt3}(1,-1,1)$. Check on basis vectors: $R_{q_1}$ is the quarter turn about $z$, so $i\mapsto j$ and $k\mapsto k$; $R_{q_2}$ is the quarter turn about $x$, so $j\mapsto k$ and $k\mapsto-j$. The composite therefore sends $i\mapsto k$ and $k\mapsto-j$. On the other hand, a rotation by $120^\circ$ about $(1,-1,1)$ permutes the three vectors $e_1,e_3,-e_2$, which make equal angles with the axis, cyclically: $e_1\mapsto e_3\mapsto-e_2\mapsto e_1$. Both descriptions agree.
\end{solution}

\begin{solution}{ex:euler4}
\idea The norm is multiplicative; expand the product.
\writeup With $p=a+bi+cj+dk$ and $q=a'+b'i+c'j+d'k$,
\[
\begin{aligned}
pq={}&(aa'-bb'-cc'-dd')\eqbr+(ab'+ba'+cd'-dc')\,i\\
&+(ac'-bd'+ca'+db')\,j\eqbr+(ad'+bc'-cb'+da')\,k,
\end{aligned}
\]
and $N(p)N(q)=N(pq)$ reads
\[
\begin{aligned}
&(a^2+b^2+c^2+d^2)\eqbr\cdot(a'^2+b'^2+c'^2+d'^2)\\
&=(aa'-bb'-cc'-dd')^2\eqbr+(ab'+ba'+cd'-dc')^2\\
&\quad+(ac'-bd'+ca'+db')^2\eqbr+(ad'+bc'-cb'+da')^2.
\end{aligned}
\]
So a product of sums of four squares is a sum of four squares, which reduces the four-square theorem to primes (Volume~2).
\end{solution}

\begin{solution}{ex:rotmatrix}
\idea Apply $v\mapsto qv\bar q$ to the basis vectors, using the multiplication table.
\writeup A direct computation  gives the images of the basis vectors, which are the columns of the matrix of $R_q$:
\[
\begin{gathered}
R_q(i)=\begin{pmatrix}a^2+b^2-c^2-d^2\\2(bc+ad)\\2(bd-ac)\end{pmatrix},\qq
R_q(j)=\begin{pmatrix}2(bc-ad)\\a^2-b^2+c^2-d^2\\2(cd+ab)\end{pmatrix},\\
R_q(k)=\begin{pmatrix}2(bd+ac)\\2(cd-ab)\\a^2-b^2-c^2+d^2\end{pmatrix}.
\end{gathered}
\]
For instance the first column is $qi\bar q$: with $q=a+x$, $qi\bar q=(a+x)i(a-x)$, and for $x=bi+cj+dk$ one finds the coefficient of $i$ to be $a^2+b^2-c^2-d^2$. The trace is $3a^2-(b^2+c^2+d^2)=3a^2-(1-a^2)=4a^2-1$. With $a=\cos\frac\theta2$, $4a^2-1=2(1+\cos\theta)-1=1+2\cos\theta$, the trace of a rotation by $\theta$ (in a basis $u,w,u\times w$ its matrix has diagonal $1,\cos\theta,\cos\theta$).
\end{solution}

\begin{solution}{ex:conjS3}
\idea Conjugation preserves the real part and rotates the pure part; $S^3$ acts transitively on directions.
\writeup Conjugation preserves real parts: $\operatorname{Re}x=\frac12(x+\bar x)$, and $\overline{pqp^{-1}}=p\bar qp^{-1}$ for unit $p$. Conversely, write unit quaternions as $q=\cos\alpha+\sin\alpha\,u$, $q'=\cos\alpha+\sin\alpha\,u'$ with the same real part $\cos\alpha$, $\alpha\in[0,\pi]$, and $u,u'$ unit pure. By Theorem~\ref{thm:su2so3} some $p$ has $R_p(u)=u'$ (there is a rotation mapping any unit vector to any other), and then $pqp^{-1}=\cos\alpha+\sin\alpha\,u'=q'$. For rotations: $R_q$ and $R_{q'}$ are conjugate in $SO(3)$ iff $q'$ is conjugate to $q$ or to $-q$ (because $R$ is onto with kernel $\pm1$), iff $\operatorname{Re}q'=\pm\operatorname{Re}q$, iff $\cos\frac{\theta'}2=\pm\cos\frac\theta2$, iff the angles in $[0,\pi]$ coincide (replacing $u$ by $-u$ changes $\theta$ into $-\theta$, so angles can be taken in $[0,\pi]$).
\end{solution}

\begin{solution}{ex:2T}
\idea A finite subset of a group closed under products is a subgroup; then count elements by order.
\writeup Let $T$ be the set of the $24$ quaternions. Its elements are exactly the quaternions of norm $1$ whose four coordinates are either all integers or all halves of odd integers: with integer coordinates, $a^2+b^2+c^2+d^2=1$ forces one coordinate $\pm1$ and the others $0$; with half-odd coordinates $\frac{m_t}2$, $\sum m_t^2=4$ with all $m_t$ odd forces every $m_t=\pm1$. In the language of Exercise~\ref{ex:hurwitz}(a), $T$ is the set of elements of norm $1$ of the Hurwitz quaternions $\mathcal H$. That part of Exercise~\ref{ex:hurwitz} shows, by a direct computation that does not use the present exercise, that $\mathcal H$ is closed under products. Since the norm is multiplicative, the product of two elements of $T$ lies in $\mathcal H$ and has norm $1$, so it lies in $T$; and $q^{-1}=\bar q$ has the same kind of coordinates as $q$. So $T$ is a subgroup of $S^3$. (For instance $i\cdot\frac12(1+i+j+k)=\frac12(-1+i-j+k)$ and $\frac12(1+i+j+k)^2=\frac12(-1+i+j+k)$, as the reader can check with Proposition~\ref{prop:quat}.) The image has $24/2=12$ elements. The orders in $SO(3)$: $\pm i,\pm j,\pm k$ give half-turns ($\cos\frac\theta2=0$), so $3$ elements of order $2$; the $16$ elements with real part $\pm\frac12$ give rotations by $120^\circ$ ($\cos\frac\theta2=\pm\frac12$), so $8$ elements of order $3$. A group of order $12$ in $SO(3)$ is $A_4$ or $D_6$ (Theorem~\ref{thm:finiteso3}), and $D_6$ has $7$ elements of order $2$; so the image is $A_4$: the rotation group of the tetrahedron with vertices $(1,1,1)$, $(1,-1,-1)$, $(-1,1,-1)$, $(-1,-1,1)$.
\end{solution}

\begin{solution}{ex:2I}
\idea Elements of order $2$ in $S^3$ are square roots of $1$, and a pure quaternion squares to a negative number.
\writeup $2I$ is the preimage of a group of order $60$ under a two-to-one surjective homomorphism, so $|2I|=120$. If $q=a+x$ with $x$ pure and $q^2=1$, then $q^2=a^2-|x|^2+2ax=1$ (Proposition~\ref{prop:quat}(d)), so $ax=0$. If $x\ne0$, then $a=0$ and $q^2=-|x|^2<0$: impossible. So $x=0$, $q=\pm1$, and $-1$ is the only element of order $2$. But $\mathfrak S_5$ has many elements of order $2$ (the $10$ transpositions, for instance), and so has $A_5\times\Z/2\Z$ (all $(a,0)$ with $a$ a double transposition, and $(e,1)$). So $2I$ is isomorphic to neither, although it has order $120$ and maps onto $A_5$. (In fact $2I\cong SL_2(\mathbb F_5)$.)
\end{solution}

\begin{solution}{ex:hurwitz}
\idea Copy the Gaussian proof: a Euclidean division gives principal (left) ideals, and a prime that divides a product without dividing the factors splits into two factors of equal norm. The only new feature is non-commutativity, which forces us to work with left ideals.
\writeup (a) $\mathcal H=\Z^4+\Z\omega$ with $\omega=\frac12(1+i+j+k)$, as an additive group. It is closed under products because the Lipschitz quaternions $\Z+\Z i+\Z j+\Z k$ are, and $\omega^2=\omega-1$, $i\omega=\frac12(-1+i-j+k)$, $\omega i=\frac12(-1+i+j-k)$ and the similar products with $j$, $k$ lie in $\mathcal H$ (all four coordinates are half-odd). Conjugation changes signs, so it preserves $\mathcal H$. For four half-odd integers $\frac{m_t}2$, $\sum\frac{m_t^2}4$ is an integer because each odd $m_t^2\equiv1\pmod8$, so $\sum m_t^2\equiv4\pmod8$. Units: norm $1$ with integer coordinates gives $\pm1,\pm i,\pm j,\pm k$; with half-odd coordinates, $\sum m_t^2=4$ forces all $m_t=\pm1$: the $16$ elements $\frac12(\pm1\pm i\pm j\pm k)$.

(b) Round each coordinate of $x$ to a nearest integer, giving $m\in\Z^4$ with all $|x_t-m_t|\le\frac12$, so $N(x-m)\le1$, with equality only if all four differences are $\pm\frac12$; but then $x\in m+\frac12(\pm1\pm i\pm j\pm k)\subseteq\mathcal H$, and $h=x$ works. Otherwise $h=m$ works. For $a,b$: let $x=ab^{-1}$, $h$ as above, $q=h$ and $r=a-qb=(x-h)b$; then $r\in\mathcal H$ and $N(r)=N(x-h)N(b)<N(b)$. If $I$ is a non-zero left ideal (stable under addition and left multiplication by $\mathcal H$), let $b\in I$ have the least non-zero norm; for $a\in I$, $r=a-qb\in I$ has $N(r)<N(b)$, so $r=0$: $I=\mathcal Hb$.

(c) Let $\alpha=1+xi+yj$, with $N(\alpha)=1+x^2+y^2=pm$. By (b), $\mathcal Hp+\mathcal H\alpha=\mathcal H\delta$. As $p\in\mathcal H\delta$, $p=\gamma\delta$, so $p^2=N(\gamma)N(\delta)$.
If $N(\delta)=1$, $\delta$ is a unit and $1=up+v\alpha$ for some $u,v\in\mathcal H$; multiplying on the right by $\bar\alpha$, $\bar\alpha=p(u\bar\alpha+vm)$ (as $p$ is central), so $\bar\alpha/p\in\mathcal H$, which is false: its real part $\frac1p$ is neither an integer nor a half-integer.
If $N(\delta)=p^2$, $\gamma$ is a unit, $\mathcal H\delta=\mathcal Hp$, and $\alpha\in\mathcal Hp$, i.e.\ $\alpha/p\in\mathcal H$: false for the same reason.
So $N(\delta)=p$. If $\delta$ has integer coordinates, $p$ is a sum of four squares. Otherwise its coordinates are half-odd; subtracting even integers, $\delta=2\mu+\varepsilon$ with $\mu\in\Z^4$ and $\varepsilon=\frac12(\pm1\pm i\pm j\pm k)$, a unit. Then $\delta\bar\varepsilon=2\mu\bar\varepsilon+1$, with $\mu\bar\varepsilon\in\mathcal H$, so $2\mu\bar\varepsilon$ has integer coordinates, and so has $\delta\bar\varepsilon$, of norm $N(\delta)=p$. Finally $2=1^2+1^2+0^2+0^2$, and Euler's identity (Exercise~\ref{ex:euler4}) extends the result to all positive integers.
\end{solution}

\begin{solution}{pb:sylvester}
\idea For (a), something must decrease: here it is the numerator of the remainder. For (b), the recursion $s_{k+1}-1=
s_k(s_k-1)$ turns $\frac1{s_k}$ into a difference, exactly as $\frac1{n(n+1)}=\frac1n-\frac1{n+1}$. For (c), a sequence that
squares at each step is $E^{2^k}$ to first order; taking logarithms turns the small correction $\frac14$ into a convergent
series.

\writeup (a) Let $n=\lceil\frac ba\rceil\ge2$, so that $\frac1n\le\frac ab<\frac1{n-1}$. The remainder is
$\frac ab-\frac1n=\frac{an-b}{bn}$, and $0\le an-b<a$ because $a(n-1)<b$. So the numerators (before reduction) decrease
strictly and the process stops. The remainder is $<\frac1{n-1}-\frac1n=\frac1{n(n-1)}\le\frac1n$, so the next unit fraction
is smaller than $\frac1n$: the fractions are distinct. For $\frac4{13}$: $n=4$, remainder $\frac3{52}$; then $n=18$,
remainder $\frac3{52}-\frac1{18}=\frac1{468}$. So $\frac4{13}=\frac14+\frac1{18}+\frac1{468}$.

(b) Since $s_{k+1}-1=s_k^2-s_k=s_k(s_k-1)$, we get $\frac1{s_k-1}-\frac1{s_{k+1}-1}=\frac{s_k-1}{s_k(s_k-1)}=\frac1{s_k}$, and the
sum telescopes: $\sum_{k=0}^N\frac1{s_k}=\frac1{s_0-1}-\frac1{s_{N+1}-1}=1-\frac1{s_{N+1}-1}$. Iterating the same identity,
$s_{k+1}-1=s_ks_{k-1}\cdots s_0(s_0-1)=s_0s_1\cdots s_k$, so $s_{k+1}\equiv1\pmod{s_j}$ for every $j\le k$: the $s_k$ are
pairwise coprime. (Their prime factors are therefore all distinct, which proves again that there are infinitely many primes,
as the Fermat numbers do in Volume~2, Chapter~2.)

(c) Let $u_k=s_k-\frac12$. Then $u_{k+1}=s_k^2-s_k+\frac12=u_k^2+\frac14$, and $u_0=\frac32$, so $(u_k)$ increases to infinity.
Let $c_k=2^{-k}\log u_k$. Then
\[
c_{k+1}-c_k=2^{-k-1}\log\Bigl(1+\frac1{4u_k^2}\Bigr)\in\Bigl]0,\frac{2^{-k-1}}{4u_k^2}\Bigr],
\]
a summable positive sequence; so $c_k$ increases to a limit $c$. Put $E=e^{c/2}$, so $E^{2^{k+1}}=e^{2^kc}$. Then
\[
2^kc-\log u_k=2^k\sum_{j\ge k}(c_{j+1}-c_j)=\sum_{j\ge k}2^{k-j-1}\log\Bigl(1+\frac1{4u_j^2}\Bigr),
\]
which lies in $\bigl]0,\frac1{4u_k^2}\bigr]$, because $u_j\ge u_k$ and $\sum_{j\ge k}2^{k-j-1}=1$. Hence
$u_k<E^{2^{k+1}}\le u_ke^{1/(4u_k^2)}$. Since $e^t-1\le te^t$ and $u_k\ge\frac32$,
\[
u_k\bigl(e^{1/(4u_k^2)}-1\bigr)\le\frac{e^{1/9}}{4u_k}\le\frac{1.12}6<\frac12.
\]
So $s_k-\frac12<E^{2^{k+1}}<s_k$, that is $s_k<E^{2^{k+1}}+\frac12<s_k+\frac12$, and $s_k=\bigl\lfloor E^{2^{k+1}}+\frac12
\bigr\rfloor$. Numerically $E=1.2640847353\ldots$
\tools Volume~2, Chapters~1--2 (coprimality, infinitely many primes); Volume~1, Chapter~4 (monotone sequences, series);
Volume~5, Chapter~13 (asymptotics of fast recursions).
\leads Aho and Sloane (1973) studied such doubly exponential recursions. Erdős and Graham asked many questions on Egyptian
fractions; one of them, that every $\frac4n$ is a sum of three unit fractions (Erdős--Straus, 1948), is still open.
\end{solution}

\begin{solution}{pb:jensen}
\idea Each quadratic factor is $|x-e^{i\varphi}|^2$, and the angles $\theta+\frac{2\pi k}n$ are the arguments of the $n$-th
roots of $e^{in\theta}$, so the product is $|x^n-e^{in\theta}|^2$. Taking logarithms gives a Riemann sum for Poisson's
integral, exactly computable. Jensen's formula is then the sum of Poisson's integrals over the roots.

\writeup (a) $x^2-2x\cos\varphi+1=(x-e^{i\varphi})(x-e^{-i\varphi})$, and the $e^{i(\theta+2\pi k/n)}$, $0\le k<n$, are the $n$
roots of $X^n=e^{in\theta}$. So $\prod_k(x-e^{i(\theta+2\pi k/n)})=x^n-e^{in\theta}$, and likewise with $-\theta$. The product
is $(x^n-e^{in\theta})(x^n-e^{-in\theta})=x^{2n}-2x^n\cos n\theta+1$.

(b) For $|x|\ne1$, $f(\theta)=\log(1-2x\cos\theta+x^2)$ is continuous (as $1-2x\cos\theta+x^2=|x-e^{i\theta}|^2>0$), so its
Riemann sums converge (Volume~1, Chapter~9). By (a) with $\theta=0$,
\[
\frac{2\pi}n\sum_{k=0}^{n-1}f\Bigl(\frac{2\pi k}n\Bigr)=\frac{2\pi}n\log(x^n-1)^2=\frac{4\pi}n\log|x^n-1|.
\]
If $|x|<1$, $\log|x^n-1|\to0$ and the limit is $0$. If $|x|>1$, $\frac{4\pi}n\log|x^n-1|=4\pi\log|x|+\frac{4\pi}n\log|1-x^{-n}|\to
4\pi\log|x|$. So $\int_0^{2\pi}f$ is $0$, or $4\pi\log|x|$, and as $f(2\pi-\theta)=f(\theta)$, $\int_0^\pi f$ is half of that.

(c) Write $z_k=r_ke^{i\alpha_k}$, $r_k\ne1$. Then $\log|e^{i\theta}-z_k|=\frac12\log(1-2r_k\cos(\theta-\alpha_k)+r_k^2)$. Over a
full period the shift by $\alpha_k$ does not matter, so by (b) (and its full-period version)
$\frac1{2\pi}\int_0^{2\pi}\log|e^{i\theta}-z_k|\dd\theta=\frac1{2\pi}\cdot\frac12\cdot4\pi\log\max(1,r_k)=\log\max(1,|z_k|)$. Add
these up, with $\log|a|$ for the constant.

The quantity $M(P)=|a|\prod_k\max(1,|z_k|)$ is the Mahler measure, and (c) says $\log M(P)$ is the average of
$\log|P|$ on the circle. If $P\in\Z[X]$ is irreducible with $P(0)\ne0$ and $M(P)=1$, then $|a|\ge1$ forces $|a|=1$ and all
$|z_k|\le1$: by Problem~\ref{pb:kronecker}, $P=\pm\Phi_n$.
\tools Volume~1, Chapters~6 and~9 (roots of unity, Riemann sums); Volume~5, Chapter~6 (integrals); Problem~\ref{pb:kronecker}.
\leads (c) is Jensen's formula (1899) on the unit circle, for polynomials. For holomorphic functions it counts zeros in a
disc, and Volume~6, Chapter~6 uses it to bound the zeros of $\zeta$. The formula also holds when $P$ has roots on the circle,
since the logarithmic singularities are integrable.
\end{solution}

\begin{solution}{pb:mersennediv}
\idea Same game as for Fermat numbers. The order of $2$ modulo a factor $q$ is the prime $p$, so $p\mid q-1$. And
$2^{p+1}=2\cdot2^p\equiv2$ exhibits $2$ as a square, which pins $q$ modulo $8$. For (b), Euler's criterion turns
``$2$ is a square modulo $2p+1$'' into ``$2^p\equiv1\pmod{2p+1}$''.

\writeup (a) $2^p\equiv1\pmod q$ and $2\not\equiv1$, so the order of $2$ is $p$, which divides $q-1$; as $q$ is odd, $2p\mid q-1$.
Moreover $p+1$ is even and $2\equiv2^{p+1}=\bigl(2^{(p+1)/2}\bigr)^2$, so $\bigl(\frac2q\bigr)=1$ and $q\equiv\pm1\pmod8$
(Volume~2, Chapter~5).

(b) If $p\equiv3\pmod4$, $q=2p+1\equiv7\pmod8$, so $\bigl(\frac2q\bigr)=1$ and Euler's criterion gives $2^{(q-1)/2}=2^p\equiv1\pmod q$. So
$q\mid M_p$, and $q<M_p$ when $p>3$, so $M_p$ is composite: $23\mid M_{11}=2047$, $47\mid M_{23}$, $167\mid M_{83}$.
\tools Volume~2, Chapters~4 and~5.
\leads Euler stated (b) in 1750. By (a), the trial divisions for $M_p$ need only the primes $2kp+1\equiv\pm1\pmod8$, which is
how the GIMPS project eliminates most exponents before running the Lucas--Lehmer test of Problem~\ref{pb:lucaslehmer}. Cole's
factor $193707721=2\cdot1445580\cdot67+1$ is of the form of (a).
\end{solution}

\begin{solution}{pb:dice}
\idea Once all six numbers are different, a rotation that fixes a numbering fixes every face, so it is the identity. The
action is free, and every orbit has $24$ elements.

\writeup (a) The rotations of the cube form a group of order $24$ (it acts transitively on the $6$ faces, and a face is fixed
by the $4$ rotations about its axis: orbit--stabiliser). It acts on the $6!=720$ numberings. A rotation that fixes a numbering
maps each face to the face with the same number, that is, to itself; the only such rotation is the identity. So all orbits
have $24$ elements, and there are $\frac{720}{24}=30$ dice.

(b) A standard numbering puts the pairs $\{1,6\}$, $\{2,5\}$, $\{3,4\}$ on the three axes ($3!$ ways) and chooses which end of
each axis gets the smaller number ($2^3$ ways): $48$ numberings, permuted by the rotations, freely as before. So there are
$\frac{48}{24}=2$ standard dice.

(c) A mirror image is $f\circ s$, where $f$ is the numbering and $s$ a reflection of the cube. If $f\circ s$ and $f$ were the
same die, then $f\circ s=f\circ g$ for some rotation $g$, so $s=g$, because $f$ is a bijection. But a reflection is not a
rotation (its determinant is $-1$). So the $30$ dice form $15$ mirror pairs, and the two standard dice are mirror images of each other.
\tools Volume~3, Chapters~1 and~3 (rotations of the cube, orbit--stabiliser).
\leads Dice are \emph{chiral}, like the molecules of Pólya's story: a molecule and its mirror image can be different
substances, as are the two forms of carvone, one smelling of spearmint and the other of caraway.
\end{solution}

\begin{solution}{pb:ksubsets}
\idea Inner products of permutation characters count orbits on pairs, and for pairs of subsets the only invariant is the size of
the intersection. So the matrix of inner products is $\min(k,l)+1$, the matrix of the ``staircase'' $\pi_k=\chi_0+\dots+\chi_k$
with orthonormal $\chi_j$. The Livingstone--Wagner theorem then falls out by restriction to $G$.

\writeup (a) $\pi_k\pi_l(\sigma)$ is the number of pairs $(A,B)$, $|A|=k$, $|B|=l$, fixed by $\sigma$, so
$\langle\pi_k,\pi_l\rangle$ is the number of orbits of $\mathfrak S_n$ on such pairs ($\pi_k$ is real-valued). If $|A\cap B|=|A'\cap B'|=i$,
the four regions $A\cap B$, $A\setminus B$, $B\setminus A$ and the complement of $A\cup B$ have the same sizes for both pairs, so a
permutation maps one pair to the other. Hence the orbits correspond to the possible values of $i$, which are
$\max(0,k+l-n),\dots,\min(k,l)$; for $k,l\le\frac n2$ that is $0,\dots,\min(k,l)$, and there are $\min(k,l)+1$.

(b) With $\pi_{-1}=0$, for $0\le k\le l\le\frac n2$:
\[
\langle\chi_k,\chi_l\rangle=\langle\pi_k,\pi_l\rangle-\langle\pi_{k-1},\pi_l\rangle-\langle\pi_k,\pi_{l-1}\rangle+\langle\pi_{k-1},\pi_{l-1}\rangle,
\]
which is $(k+1)-k-(k+1)+k=0$ if $k<l$ and $(k+1)-k-k+k=1$ if $k=l$ (for $k=0$ drop the terms with $\pi_{-1}$; the result is the
same). Each $\chi_k=\sum_im_i\psi_i$ is an integer combination of irreducible characters $\psi_i$ with $\sum m_i^2=1$, so
$\chi_k=\pm\psi$ for one irreducible $\psi$; and $\chi_k(e)=\binom nk-\binom n{k-1}>0$ for $k\le\frac n2$, so $\chi_k=\psi$. They are
distinct because they are orthogonal.

(c) The number of orbits of $G$ on $k$-subsets is $\frac1{|G|}\sum_{g\in G}\pi_k(g)=\langle\pi_k|_G,1\rangle_G$ by Burnside. By (b),
$\pi_k=\chi_0+\chi_1+\dots+\chi_k$, so this number is $\sum_{j\le k}\langle\chi_j|_G,1\rangle_G$. Each $\chi_j$ is the character of a
representation of $\mathfrak S_n$, hence $\chi_j|_G$ of a representation of $G$, and $\langle\chi_j|_G,1\rangle_G$ is the
dimension of its fixed vectors, at least $0$. So the sum is non-decreasing in $k$.
\tools Volume~3, Chapters~3 and~11 (Burnside, orthogonality of characters); Volume~1, Chapter~2 (binomial coefficients).
\leads The $\chi_k$ are the characters of the partitions $(n-k,k)$ in Young's theory. For $n=4$ they have degrees $1,3,2$; for
$n=5$, $1,4,5$. Livingstone and Wagner (1965) proved (c) to show that a group that is transitive on $k$-subsets is also transitive
on $(k-1)$-subsets.
\end{solution}

\begin{solution}{pb:egz}
\idea Chevalley--Warning turns ``an equation of low degree has a non-trivial zero'' into a congruence. The polynomial $x^{p-1}$
is the indicator of $x\ne0$, so the two equations say ``choose $p$ indices whose $a_i$ add to $0$''.

\writeup (a) For $y\in\mathbb F_p$, $1-y^{p-1}$ is $1$ if $y=0$ and $0$ otherwise (Fermat). So the number $Z$ of common zeros
satisfies, in $\mathbb F_p$,
\[
Z\equiv\sum_{x\in\mathbb F_p^m}\prod_{j=1}^r\bigl(1-f_j(x)^{p-1}\bigr).
\]
The product is a polynomial of degree at most $(p-1)\sum_j\deg f_j<(p-1)m$, so each of its monomials $x_1^{e_1}\cdots x_m^{e_m}$ has
some $e_i<p-1$. Its sum over $\mathbb F_p^m$ is $\prod_i\sum_{x_i}x_i^{e_i}$, and $\sum_{x\in\mathbb F_p}x^e=0$ for $0\le e<p-1$
(Exercise~9.6, with $0^0=1$). So $Z\equiv0\pmod p$.

(b) Let $a_1,\dots,a_{2p-1}$ be the integers, read in $\mathbb F_p$. The system $f_1=f_2=0$ in $m=2p-1$ variables has degrees
$2(p-1)<m$, and $x=0$ is a common zero; by (a) there is another, $x\ne0$. Let $I=\{i:x_i\ne0\}$. Then $x_i^{p-1}=1$ for $i\in I$ and
$0$ otherwise, so $f_2(x)=|I|\equiv0$ and $f_1(x)=\sum_{i\in I}a_i\equiv0$. Since $0<|I|\le2p-1$, $|I|=p$.

(c) Take $2mn-1$ integers. As long as at least $2n-1$ remain, the theorem for $n$ lets us remove $n$ of them with sum divisible by
$n$. After $j$ removals $2mn-1-jn$ remain, which is at least $2n-1$ for $j\le2m-2$; so we get $2m-1$ disjoint groups, with sums
$nb_1,\dots,nb_{2m-1}$. The theorem for $m$ gives $m$ of the $b_i$ with sum divisible by $m$, and the union of these $m$ groups is a
set of $mn$ integers with sum divisible by $mn$. Since the theorem holds for $1$ and for primes, it holds for all $n$.

(d) Take $n-1$ zeros and $n-1$ ones: $n$ of them contain between $1$ and $n-1$ ones, so their sum is not divisible by $n$.

\tools Exercises~9.6 and~9.8 (power sums, Chevalley--Warning for one polynomial); Volume~2, Chapter~3 (Fermat).
\leads The original 1961 proof used Cauchy--Davenport instead. In the plane, Reiher proved in 2007 Kemnitz's conjecture: among
any $4n-3$ lattice points there are $n$ whose centroid is a lattice point.
\end{solution}

\begin{solution}{pb:abelian}
\idea The structure theorem reduces abelian groups of order $p^e$ to partitions of $e$, and the generating function of partitions
becomes an Euler product of zeta values. The average is governed by the simple pole of $\zeta(s)$ at $1$; the other factors are a
convergent correction.

\writeup (a) A finite abelian group is the direct product of its Sylow subgroups, and an abelian $p$-group of order $p^e$ is
$\prod_i\Z/p^{\lambda_i}\Z$ for a unique partition $\lambda$ of $e$ (Corollary~7.7). So $A$ is multiplicative with $A(p^e)=p(e)$. Its
Euler factor is $\sum_ep(e)p^{-es}=\prod_{k\ge1}(1-p^{-ks})^{-1}$, and the product over $p$ is $\prod_k\zeta(ks)$ (all products converge
absolutely for $s>1$).

(b) Write $\prod_k\zeta(ks)=\zeta(s)G(s)$ with $G(s)=\prod_{k\ge2}\zeta(ks)=\sum_mg(m)m^{-s}$. Here $g\ge0$, and $g$ is supported on the
powerful numbers (its Euler factor at $p$ has no term $p^{-s}$). Then $A=\mathbf1*g$ and
\[
\sum_{n\le x}A(n)=\sum_{m\le x}g(m)\Bigl\lfloor\frac xm\Bigr\rfloor=x\sum_mg(m)m^{-1}-x\sum_{m>x}\frac{g(m)}m+O\Bigl(\sum_{m\le x}g(m)\Bigr).
\]
Write $G(s)=\zeta(2s)K(s)$ with $K(s)=\prod_{k\ge3}\zeta(ks)=\sum_jh(j)j^{-s}$, $h\ge0$, convergent at every $s>\frac13$. Then
$\sum_{m\le y}g(m)=\sum_jh(j)\lfloor\sqrt{y/j}\rfloor\le\sqrt y\sum_jh(j)j^{-1/2}=K(\tfrac12)\sqrt y$. Writing $G(y)=\sum_{m\le y}g(m)$, partial summation gives
$\sum_{m>x}\frac{g(m)}m=-\frac{G(x)}x+\int_x^\infty\frac{G(y)}{y^2}dy\le\int_x^\infty\frac{K(\frac12)}{y^{3/2}}dy=\frac{2K(\frac12)}{\sqrt x}$. So the sum is $G(1)x+O(\sqrt x)$, with
$G(1)=\prod_{k\ge2}\zeta(k)=2.2948\ldots$

\tools Corollary~7.7 (finitely generated abelian groups); Volume~5, Proposition~2.11 (generating function of $p(n)$);
Volume~6, Chapter~2 (convolutions and partial summation).
\leads Erdős and Szekeres proved (b) in 1934. The truth is $\sum_{n\le x}A(n)=c_1x+c_2x^{1/2}+c_3x^{1/3}+o(x^{1/3})$ with $c_k=\prod_{j\ne k}\zeta(\frac jk)$, so
$c_2<0$ because $\zeta(\frac12)<0$; the best exponent in the error term is an open problem. The individual values of $A(n)$ can be large: $A(2^{e})=p(e)$ grows like
$e^{c\sqrt e}$.
\end{solution}

\begin{solution}{pb:dodecaturns}
\idea A law that is constant on conjugacy classes is diagonalised by the characters: each irreducible representation turns convolution into
multiplication by a number. For $A_5$ and $3$-cycles the numbers are $1,0,0,\frac14,-\frac15$.

\writeup (a) Let $\rho$ be irreducible of dimension $d$ and $A=\sum_hf(h)\rho(h)$ with $f$ a class function. Then $\rho(g)A\rho(g)^{-1}=\sum_hf(h)\rho(ghg^{-1})=A$ after the
substitution $h\mapsto g^{-1}hg$, so $A$ is a scalar $cI$ by Schur's lemma (Lemma~11.4); taking traces, $cd=\hat f(\chi)$. For two functions, with
$(f*g)(x)=\sum_hf(h)g(h^{-1}x)$, one checks $\sum_x(f*g)(x)\rho(x)=\bigl(\sum_hf(h)\rho(h)\bigr)\bigl(\sum_yg(y)\rho(y)\bigr)$. The law of $X_t=\xi_1\cdots\xi_t$ is $\mu^{*t}$, which is again a
class function, so $\widehat{\mu^{*t}}(\chi)=d\,(\hat\mu(\chi)/d)^t$.

The irreducible characters form an orthonormal basis of the class functions (Theorems~11.8 and~11.10). For a class function $F$,
$\langle F,\psi\rangle=\frac1{|G|}\sum_hF(h)\overline{\psi(h)}=\frac1{|G|}\hat F(\bar\psi)$, and $\bar\psi$ is again an irreducible character. Hence
$F=\frac1{|G|}\sum_\chi\hat F(\chi)\bar\chi$, and with $F=\mu^{*t}$ this is the formula.

(b) The character table of $A_5$, on the classes of $e$, $(12)(34)$, $(123)$, $(12345)$, $(13524)$, has rows $(1,1,1,1,1)$, $(3,-1,0,\varphi,1-\varphi)$,
$(3,-1,0,1-\varphi,\varphi)$, $(4,0,1,-1,-1)$, $(5,1,-1,0,0)$, with $\varphi=\frac{1+\sqrt5}2$. For $\mu$ uniform on the $3$-cycles, $\frac{\hat\mu(\chi)}{\chi(1)}=\frac{\chi((123))}{\chi(1)}$ takes the values $1,0,0,\frac14,-\frac15$. For $g=e$
and $t\ge1$ the formula gives
\[
\mu^{*t}(e)=\frac1{60}\Bigl(1+16\cdot4^{-t}+25\Bigl(-\frac15\Bigr)^t\Bigr).
\]
For $t=1$ this is $0$, as it should be, and for $t=2$ it is $\frac3{60}$: of the $400$ pairs of turns, $20$ undo each other.

(c) The rotations by $\pm\frac{2\pi}5$ about the six vertex axes form one conjugacy class of $12$ elements of order $5$. At such a $g$ the characters of
dimensions $4$ and $5$ are $-1$ and $0$, and the $3$-dimensional ones have ratio $0$, so $\mu^{*t}(g)=\frac1{60}(1-4\cdot4^{-t})$ and the probability asked for is
$\frac15(1-4^{1-t})$, for instance $\frac3{20}$ after two turns. In general $|\mu^{*t}(g)-\frac1{60}|\le\frac1{60}(16\cdot4^{-t}+25\cdot5^{-t})\le\frac{41}{60}4^{-t}$ for every $g$, so the distance to the uniform
law is at most $\frac{41}2\,4^{-t}$; it is not smaller order, since at $e$ the deviation is about $\frac{16}{60}4^{-t}$.

\tools{Lemma~11.4, Theorems~11.8 and 11.10, Exercise~11.7 (the table of $A_5$), Theorem~12.5.}
\leads Diaconis and Shahshahani used exactly this method in 1981, with the characters of $\mathfrak S_n$, to show that random transpositions mix a deck
after about $\frac12n\ln n$ steps and not before, one of the first cutoff results. Volume~4, Chapter~8 proves the upper bound by Broder's
probabilistic argument instead.
\end{solution}

\section*{Sources and further reading}

\begin{itemize}
\item W.~R.~Hamilton, \emph{On quaternions, or on a new system of imaginaries in algebra}, Philosophical Magazine 25 (1844), 10--13, and the following instalments. The discovery, a few months after the walk.
\item F.~Klein, \emph{Vorlesungen über das Ikosaeder und die Auflösung der Gleichungen vom fünften Grade}, Teubner, 1884 (English: \emph{Lectures on the Icosahedron}, Dover). Finite rotation groups, and the quintic solved with the icosahedron.
\item A.~Hurwitz, \emph{Über die Zahlentheorie der Quaternionen}, Nachr.\ Ges.\ Wiss.\ Göttingen (1896), 313--340. The ring $\mathcal H$ and the four-square theorem.
\item J.~Stillwell, \emph{Naive Lie Theory}, Springer, 2008. Quaternions, $SU(2)$, $SO(3)$, and the first steps of Lie theory.
\item J.~H.~Conway and D.~A.~Smith, \emph{On Quaternions and Octonions}, A K Peters, 2003. The finite groups of quaternions, Hurwitz arithmetic, and the octonions.
\item M.~A.~Armstrong, \emph{Groups and Symmetry}, Springer, 1988, Chapter~19. The classification of finite rotation groups by counting poles.
\item K.~Shoemake, \emph{Animating rotation with quaternion curves}, SIGGRAPH Computer Graphics 19 (1985), 245--254. Quaternion interpolation (slerp) in computer graphics.
\end{itemize}
